\documentclass[11pt,a4paper]{article}
\usepackage[T1]{fontenc}
\usepackage{lmodern,amsmath,amssymb,amsthm,mathtools,booktabs,microtype}
\usepackage[margin=26mm]{geometry}
\usepackage[colorlinks=true,linkcolor=blue,urlcolor=blue]{hyperref}
\setlength{\emergencystretch}{2em}
\newtheorem{theorem}{Theorem}[section]
\newtheorem{lemma}[theorem]{Lemma}
\newtheorem{corollary}[theorem]{Corollary}
\theoremstyle{remark}\newtheorem{remark}[theorem]{Remark}
\newcommand{\Q}{\mathbb Q}
\DeclareMathOperator{\Tr}{Tr}
\title{Quadratic rescaling of elliptic reference and height coordinates}
\author{}\date{16 November 2026}
\begin{document}\maketitle
\begin{abstract}
Two quadrics in three field coordinates describe an elliptic curve with
two points removed. The reference reciprocal abscissas are linear in
these coordinates, while the two logarithmic ratios have a common
denominator. A homogeneous rescaling preserves quadratic equations for
two numerator points, their difference, and both ratios. Over a quadratic
number field, the original rational trace of the difference is recovered
with one additional rational coordinate. The reference, point, ratio,
and height blocks use 32 rational coordinates on the nonzero-scale
locus. An explicit quadratic guard adds one rational coordinate,
giving the same full block count of 33 as the unscaled source.
The statement concerns these blocks and does not supply a detector
conversion or a smaller complete defining polynomial.
\end{abstract}

\section{A reference chart with two points removed}
Let $k$ have characteristic different from two. Fix
\[
 E:y^2=x^3-Dx,\qquad P=(a,b),\qquad T_0=(0,0),
\]
with $abD\ne0$ and $b^2=a^3-Da$. Put
\[
 c=b/a,\qquad\delta=D/a=a-c^2.
\]
Consider the affine curve $\mathcal C$ with coordinates
$(\ell,J,\alpha)$ and equations
\begin{equation}\label{eq:compact}
 J=\ell^2,\qquad
 \alpha(\alpha-\delta J-1)+c\ell(1-\delta J)+2\delta J=0.
\end{equation}
Write $\gamma=\alpha-\delta J-1$.

\begin{theorem}\label{thm:chart}
The curve $\mathcal C$ is isomorphic over $k$ to
\[
                       E\setminus\{T_0,P+T_0\}.
\]
Its locus $\ell\ne0$ corresponds exactly to
$E\setminus\{O,T_0,P,P+T_0\}$. On that locus,
\begin{equation}\label{eq:ratios}
       S=\frac{y(H)}{x(H)}=\frac{\alpha}{\ell},\qquad
       T_H=\frac{y(H-P)}{x(H-P)}=\frac{\gamma}{\ell}.
\end{equation}
The reciprocal abscissas extend regularly to all of $\mathcal C$ as
the linear functions
\begin{equation}\label{eq:recips}
 f_H=\frac{1-\alpha-c\ell}{a},\qquad
 g_H=\frac{\alpha-\delta J-c\ell}{a}.
\end{equation}
The two points with $\ell=0$ are
\[
 (\ell,J,\alpha)=(0,0,1)\longleftrightarrow O,
 \qquad
 (\ell,J,\alpha)=(0,0,0)\longleftrightarrow P.
\]
Here the reciprocal abscissa at $O$ has its regular value zero.
\end{theorem}

\begin{proof}
Use the complete chart of~\cite{Three}, whose coordinates $(S,T_H,v)$
satisfy
\begin{equation}\label{eq:old}
 v(v+T_H-S)+\delta=0,\qquad
 ST_H+c(T_H-S)+2cv+2\delta=0.
\end{equation}
It is isomorphic to the stated four-point complement, and $v$ is a
unit there. Set
\[
 \ell=v^{-1},\qquad J=\ell^2,\qquad\alpha=\ell S.
\]
The first equation gives $\ell T_H=\alpha-\delta J-1=\gamma$.
Multiplication of the second equation by $\ell^2$ gives
the second equation of~\eqref{eq:compact}. Conversely, on $\ell\ne0$,
the substitutions $v=\ell^{-1}$ and~\eqref{eq:ratios} recover
both equations~\eqref{eq:old}. These are mutually inverse regular
maps. The old recovery formulas
\[
 f_H=\frac{v-S-c}{av},\qquad
 g_H=\frac{v+T_H-c}{av}
\]
give~\eqref{eq:recips}.

For the extension across $\ell=0$, the reciprocal ordinates
$h_H=-y(H)/x(H)^2$ and $k_H=-y(H-P)/x(H-P)^2$ have polynomial
expressions
\begin{equation}\label{eq:ordinates}
 \begin{split}
 h_H&=\frac{(c+\delta\ell)\alpha+c\delta J-2\delta\ell-c}{a},\\
 k_H&=\frac{(c+\delta\ell)\gamma+c-c\delta J+2\delta\ell}{a}.
 \end{split}
\end{equation}
Indeed $\ell h_H=-\alpha f_H$ and
$\ell k_H=-\gamma g_H$ follow directly from~\eqref{eq:compact}.
The same equations give
\[
 h_H^2=f_H-Df_H^3,\qquad k_H^2=g_H-Dg_H^3.
\]
The affine reciprocal cubic $h^2=f-Df^3$ is the original curve with
$T_0$ removed: its point $(f,h)=(0,0)$ represents $O$, while its
point at infinity represents $T_0$. At $\ell=0$, the equations yield
$J=0$ and $\alpha\in\{0,1\}$. Formulas~\eqref{eq:recips}
and~\eqref{eq:ordinates} then give respectively
\[
 (f_H,h_H;g_H,k_H)=(0,0;1/a,c/a)
 \quad\hbox{or}\quad(1/a,-c/a;0,0).
\]
These are exactly $(H,H-P)=(O,-P)$ and $(P,O)$.

After eliminating $J$, the second equation of~\eqref{eq:compact}
is monic of degree two in $\alpha$. Its derivative with respect to
$\alpha$ is $2\alpha-1$ at $\ell=0$, so both added points are smooth.
No component can be supported on $\ell=0$, since that fiber is
zero-dimensional and the defining polynomial is not divisible by
$\ell$. The open curve already identified is integral, hence so
is $\mathcal C$. The old functions $\ell$ and $\ell S$ extend
regularly at $O,P$ with the displayed values: $\ell$ has simple
zeros there, canceling the relevant simple pole of $S$ or $T_H$.
They have poles only at the two removed points. Thus the inverse
map also extends, and the inverse identities on the dense open
locus hold everywhere. This proves the theorem.
\end{proof}

\section{Quadratic rescaling of the point and ratio rows}
For the remaining statements assume that $k$ has characteristic zero.
Adjoin one field coordinate $\kappa$ with
\begin{equation}\label{eq:kappa}
                         \kappa=J^2.
\end{equation}
Throughout this section impose the explicit domain condition
\begin{equation}\label{eq:open}
                              \ell\ne0.
\end{equation}
For a point $U$ write its reciprocal coordinates as
\[
 f=1/x(U),\qquad h=-y(U)/x(U)^2,
 \qquad h^2=f-Df^3,
\]
with $(f,h)=(0,0)$ representing $U=O$. Use $(g,k)$ for a second
point $V$. The complete paired source~\cite{Trace} retains
$f,h,g,k,s,F_U,F_V$ and imposes
\begin{equation}\label{eq:unscaledpair}
 \begin{gathered}
 F_U=f^2,\quad F_V=g^2,\quad h^2=f-DfF_U,\\
 s(f-g)=h+k,\qquad
 s(h-k)=1-D(F_U+fg+F_V).
 \end{gathered}
\end{equation}
These rows give $U\ne V$, a finite difference $A=U-V$, and
\begin{equation}\label{eq:xA}
                         x(A)=s^2+D(f+g).
\end{equation}
They include identity numerator points and the tangent case $f=g$.

Set
\begin{equation}\label{eq:scalings}
 \begin{gathered}
 f_*=\ell^2f,\quad g_*=\ell^2g,\quad
 h_*=\ell^3h,\quad k_*=\ell^3k,\quad s_*=\ell s,\\
 F_*=\ell^4F_U,\qquad G_*=\ell^4F_V.
 \end{gathered}
\end{equation}

\begin{theorem}\label{thm:scaled}
Under~\eqref{eq:compact}, \eqref{eq:kappa} and~\eqref{eq:open},
the change~\eqref{eq:scalings} gives an exact equivalence between
\eqref{eq:unscaledpair} and the five field quadrics
\begin{equation}\label{eq:scaledpair}
 \begin{gathered}
 F_*=f_*^2,\quad G_*=g_*^2,\quad
 h_*^2=\kappa f_*-Df_*F_*,\\
 s_*(f_*-g_*)=h_*+k_*,\qquad
 s_*(h_*-k_*)=\kappa-D(F_*+f_*g_*+G_*).
 \end{gathered}
\end{equation}
The four original ratio rows are equivalent to the four field quadrics
\begin{equation}\label{eq:scaledratios}
 \begin{aligned}
 h_*L+f_*\alpha&=0,&
 (\kappa-DF_*)L+h_*\alpha&=0,\\
 k_*L_0+g_*\gamma&=0,&
 (\kappa-DG_*)L_0+k_*\gamma&=0.
 \end{aligned}
\end{equation}
The difference satisfies the exact identity
\begin{equation}\label{eq:scaledheight}
                Jx(A)=s_*^2+D(f_*+g_*).
\end{equation}
Every displayed row has total degree at most two in the retained
coordinates.
\end{theorem}

\begin{proof}
The first two equations of~\eqref{eq:scaledpair} are respectively
$\ell^4$ times the original square rows. The remaining three rows
are respectively $\ell^6$, $\ell^3$ and $\ell^4$ times their originals.
For the ratios, the unscaled rows are
\[
 hL+fS=0,\quad(1-DF_U)L+hS=0,
 \quad kL_0+gT_H=0,\quad(1-DF_V)L_0+kT_H=0.
\]
Equations~\eqref{eq:ratios} show that their rescaled versions are
$\ell^3,\ell^4,\ell^3,\ell^4$ times these rows, respectively.
Finally~\eqref{eq:scaledheight} is $\ell^2$ times~\eqref{eq:xA}.
All multipliers are nonzero, and~\eqref{eq:scalings} is invertible
on the stated locus. Thus every implication is reversible, including
all exceptional point strata admitted by~\eqref{eq:unscaledpair}.
\end{proof}

\section{Recovering the original rational trace}
Now take $k=\Q(\rho)$ with $\rho^2=2$. A field coordinate means two
rational coordinates. Keep the original rational height coordinate
$T=\Tr x(A)$ and adjoin one rational coordinate $Z$. Impose the
single field quadratic equation
\begin{equation}\label{eq:trace}
              J\left(\frac T2+\rho Z\right)
                   =s_*^2+D(f_*+g_*).
\end{equation}

\begin{corollary}\label{cor:trace}
On $\ell\ne0$, equation~\eqref{eq:trace} is equivalent to
$T=\Tr x(A)$, with $Z$ uniquely equal to the coefficient of $\rho$
in $x(A)$. In particular the recovered height coordinate depends
only on the difference $A$, as in the unscaled source.
\end{corollary}
\begin{proof}
Since $J=\ell^2\ne0$, comparison with~\eqref{eq:scaledheight}
gives $x(A)=T/2+\rho Z$. This is exactly the unique decomposition
of a field element whose trace is $T$.
\end{proof}

\begin{corollary}[An exact rational nonzero guard]\label{cor:guard}
Write $\kappa=\kappa_0+\rho\kappa_1$ in the fixed rational basis.
One rational coordinate $\zeta$ and the quadratic row
\begin{equation}\label{eq:guard}
                            \kappa_0\zeta=1
\end{equation}
enforce $\ell\ne0$ exactly. For every nonzero $\ell$ there is a
unique rational witness $\zeta$.
\end{corollary}
\begin{proof}
The equations give $\kappa=\ell^4$. If
$\ell=\ell_0+\rho\ell_1$, then
\[
 \kappa_0=\ell_0^4+12\ell_0^2\ell_1^2+4\ell_1^4.
\]
This is positive for every nonzero pair of rational coordinates
and vanishes at the zero pair. Thus $\kappa_0\ne0$ is equivalent
to $\ell\ne0$, and~\eqref{eq:guard} has exactly the witness
$\zeta=1/\kappa_0$ on that locus.
\end{proof}

For the usual ratios and comparison squares, retain
\[
 L,L_0,C_2,I_2\in k,\qquad
 C_2=(L-\overline{L_0})^2,\qquad I_2=(L-\beta)^2,
\]
where $\beta$ is an affine rational expression in retained coordinates.
These rows do not change. The coordinate comparison is therefore
\[
\begin{array}{l|r|r}
 \text{block}&\text{unscaled rational coordinates}&
                   \text{scaled rational coordinates}\\\hline
 \text{reference}&10&8\\
 \text{numerator pair and slope}&14&14\\
 \text{ratios and comparison squares}&8&8\\
 \text{trace recovery}&1&2\\\hline
 \text{subtotal on the nonzero locus}&33&32\\
 \text{additional nonzero-scale guard}&0&1\\\hline
 \text{guarded total}&33&33
\end{array}
\]
The unscaled reference block consists of the five field coordinates
of the complete reference source~\cite{Reference}; the scaled block
is $(\ell,J,\alpha,\kappa)$.

\section{Domain and detector scope}
The coordinate comparison is an equivalence of the specified algebraic
blocks on $\ell\ne0$, with~\eqref{eq:guard} providing an exact
quadratic enforcement whose cost is included above.
The equations~\eqref{eq:compact} and
\eqref{eq:kappa} themselves allow $\ell=0$. At that fiber the
scaled cubic becomes $h_*^2=-Df_*^3$, and the inverse rescaling
does not exist. Omitting the guard from a complete source therefore
requires a separate proof that its other retained equations exclude
or correctly handle every added boundary point.

Nor does the coordinate comparison automatically preserve detector
polynomials. For example the old stored square $F_U=f^2$ becomes
$F_*/\kappa$ on the nonzero locus. Clearing such denominators can
alter the required degree and the contact valuation estimates.
Theorems~\ref{thm:chart} and~\ref{thm:scaled} and
Corollaries~\ref{cor:trace} and~\ref{cor:guard} make no assertion
that those further costs vanish. The guarded block has no net
coordinate saving, and no smaller complete defining quartic follows
from this block equivalence alone.

The exact verifier in
\nolinkurl{computation/elliptic_charts/scaled_reference_v43}
checks the generic chart identities, both boundary values, every
rescaling exponent, rational splitting of~\eqref{eq:trace}, the
positive fourth-power expression in the guard, and the stated
coordinate counts. The infinite-locus assertions follow
from the proofs above and the complete cited paired source.

\begin{thebibliography}{9}
\bibitem{Three}\emph{A three-coordinate quadratic chart for an elliptic
point and its translate}, project manuscript, 15 November 2026,
\nolinkurl{papers/elliptic_charts/three_coordinate_reference_20261115_v40.tex}.
\bibitem{Trace}\emph{Trace heights on quadratic twists with $j=1728$},
project manuscript, 13 November 2026,
\nolinkurl{papers/trace_heights/trace_on_all_points_20261113_v25.tex}.
\bibitem{Reference}\emph{A complete compressed reference-point chart},
project manuscript, 14 November 2026,
\nolinkurl{papers/elliptic_charts/compressed_reference_20261114_v27.tex}.
\end{thebibliography}
\end{document}
